% !TeX root = ../../Book4.tex
% 纲要标题：### 18.2 屋顶与分式
% 纲要标签：b4:sec:1702
% 纲要位置：计划纲要.md，第 3690 行。

\section{屋顶与分式}
\label{b4:sec:1702}

导出态射可以先沿拟同构更换源对象，再作复形映射。屋顶的共同细化既给出相等判据，也使复合、加法与三角结构可以严格构造。


% 纲要内容（保留纲要核验项目）：
%
% * 屋顶图表示态射
% * 适当的分式演算
% * 同伦范畴到导出范畴
%

\subsection{态射的屋顶图表示}
\label{b4:subsec:1702:01}

\begin{definition}\label{b4:def:derived-roof}
设 \(\mathcal A\) 为局部小交换范畴，在固定的集合大小约定下记局部化函子为 \(Q:K(\mathcal A)\to D(\mathcal A)\)。\(K(\mathcal A)\) 中的图式
\[
 X\xleftarrow{s}U\xrightarrow{f}Y,\qquad s\;\text{为拟同构},
\]
称为一个\term[wuding]{屋顶}[roof]，表示 \(Q(f)Q(s)^{-1}:X\to Y\)。给定另一个屋顶 \(X\xleftarrow{s'}U'\xrightarrow{f'}Y\)，若存在对象 \(V\) 及拟同构 \(a:V\to U\)、\(b:V\to U'\)，使 \(sa=s'b\)、\(fa=f'b\)，则称两个屋顶有共同细化；所有等式都在同伦范畴中理解。
\end{definition}

共同细化把两个屋顶放在同一张图中：
\[
\begin{tikzcd}[column sep=3.2em,row sep=2.2em]
&U\arrow[dl,"s\;(\sim)"']\arrow[dr,"f"]&\\
X&V\arrow[u,"a\;(\sim)"]\arrow[d,"b\;(\sim)"']&Y\\
&U'\arrow[ul,"s'\;(\sim)"]\arrow[ur,"f'"']&
\end{tikzcd}
\]
左侧要求 \(sa=s'b\)，右侧要求 \(fa=f'b\)。这里 \(a,b\) 都是拟同构；也可以先要求共同左分母 \(sa=s'b\) 为拟同构，再由二取三性质推出这一点。符号 \(Q(s)^{-1}\) 是局部化以后得到的逆，图式本身只使用同伦范畴中实际存在的箭头，没有要求在那里选择 \(s\) 的逆映射。

\ProseParagraphBreak
在模消解的典型例子中，\(P\to M\) 是左分母，而 \(P\to N[n]\) 记录一个 Ext 余圈。屋顶因此不是额外的画图规则，而是“先换成可计算的模型，再作映射”的正式表达。接下来证明共同细化足以判别相等。

\subsection{分式演算}
\label{b4:subsec:1702:02}

屋顶的复合需要把指向同一对象的两张箭头通分；比较不同的通分选择，则需要消去两张箭头的差。下面第一项构造等式 \(ft=sg\)，使复合所需的方块交换；第二项从 \(sf=sg\) 出发，允许进一步前复合一个拟同构后得到 \(ft=gt\)，用来保证计算不依赖选择。第三、第四项分别是后复合方向的通分和消去，它们给出分母在右的分式演算。

\begin{lemma}\label{b4:lem:quasi-ore}
设 \(\mathcal A\) 为局部小交换范畴。\(K(\mathcal A)\) 中的拟同构包含恒等态射，对复合和平移封闭，并满足二取三性质。它们还满足以下四个条件。
\begin{enumerate}
\item 给定 \(f:X\to Y\) 和拟同构 \(s:Z\to Y\)，存在拟同构 \(t:W\to X\) 和 \(g:W\to Z\)，使 \(ft=sg\)。
\item 若拟同构 \(s:Y\to Y'\) 满足 \(sf=sg\)，其中 \(f,g:X\to Y\)，则存在拟同构 \(t:X'\to X\)，使 \(ft=gt\)。
\item 给定 \(f:X\to Y\) 和拟同构 \(s:X\to Z\)，存在拟同构 \(t:Y\to W\) 和 \(g:Z\to W\)，使 \(tf=gs\)。
\item 若拟同构 \(s:X'\to X\) 满足 \(fs=gs\)，其中 \(f,g:X\to Y\)，则存在拟同构 \(t:Y\to Y'\)，使 \(tf=tg\)。
\end{enumerate}
\end{lemma}
\begin{proof}
恒等态射、复合、平移及二取三性质的结论直接对各次上同调核对。对第一项，补全 \(Z\xrightarrow{s}Y\to C\to Z[1]\)。因 \(s\) 为拟同构，\(C\) 零调。再把复合 \(X\xrightarrow{f}Y\to C\) 补成旋转后的好三角 \(W\xrightarrow{t}X\to C\to W[1]\)。长正合列说明 \(t\) 是拟同构；对这两个三角使用 (TR3)，即可得到 \(g\) 及 \(ft=sg\)。

对第二项令 \(h=f-g\)。补全 \(s\) 的锥三角，记 \(C=\operatorname{Cone}(s)\)。由 \(sh=0\) 和\cref{b4:lem:pretriangle-hom}的协变 Hom 正合性，\(h\) 经 \(C[-1]\) 分解为 \(X\xrightarrow{a}C[-1]\to Y\)。这里 \(C[-1]\) 虽然零调，却未必是 \(K(\mathcal A)\) 中的零对象，不能仅由这一分解断言 \(h=0\)。把 \(a\) 补成好三角
\[
 X'\xrightarrow{t}X\xrightarrow{a}C[-1]\longrightarrow X'[1].
\]
由\cref{b4:prop:cohomology-on-homotopy}产生的上同调长正合列及 \(C[-1]\) 零调，\(t\) 是拟同构。相邻复合为零给出 \(at=0\)，所以 \(ht=0\)，即 \(ft=gt\)。零调锥在这里提供的是进一步拟同构后的消去，而非原同伦范畴中的直接消去。

对第三项，补全 \(X\xrightarrow{s}Z\to C\to X[1]\)，其中 \(C\) 零调。逆向旋转后记首箭头为 \(u:C[-1]\to X\)，再把 \(fu:C[-1]\to Y\) 补成好三角
\[
 C[-1]\xrightarrow{fu}Y\xrightarrow{t}W\to C.
\]
对逆向旋转的三角和这个三角应用 (TR3)，前两个分量取 \(1_{C[-1]},f\)，得到 \(g:Z\to W\)，并有 \(tf=gs\)。长正合列及 \(C\) 零调说明 \(t\) 是拟同构。

对第四项令 \(h=f-g\)，补全 \(X'\xrightarrow{s}X\xrightarrow{q}C\to X'[1]\)。由 \(hs=0\) 和反变 Hom 正合性，存在 \(b:C\to Y\) 使 \(h=bq\)。把 \(b\) 补成好三角 \(C\xrightarrow{b}Y\xrightarrow{t}Y'\to C[1]\)。\(C\) 零调，故 \(t\) 是拟同构，而 \(th=tbq=0\)。
\end{proof}

\begin{theorem}\label{b4:thm:roof-calculus}
设 \(\mathcal A\) 为局部小交换范畴，并采用\subsecref{b4:subsec:1701:03}的集合大小约定，记局部化函子为 \(Q:K(\mathcal A)\to D(\mathcal A)\)。\(D(\mathcal A)\) 中每个态射可由屋顶表示，两个屋顶表示同一态射当且仅当它们有共同细化。特别地，对 \(K(\mathcal A)\) 中的 \(f,g:X\to Y\)，\(Q(f)=Q(g)\) 当且仅当存在拟同构 \(t:X'\to X\) 使 \(ft=gt\)；也可用后复合的拟同构 \(Y\to Y'\) 检验。
\end{theorem}
\begin{proof}
先用共同细化定义屋顶的等价。自反、对称显然。对传递性，把两次细化到共同中间屋顶的分母用\cref{b4:lem:quasi-ore}第一项再作共同细化，二取三保证新添的箭头仍为拟同构；复合之后，左右两条边同时相等。这证明它确为等价关系。

复合 \((s,f):X\to Y\) 与 \((t,g):Y\to Z\) 时，对 \(f:U\to Y\) 和 \(t:V\to Y\) 取 \(a:W\to U\)、\(b:W\to V\)，其中 \(a\) 拟同构且 \(fa=tb\)。拼接图为
\[
\begin{tikzcd}[column sep=1.8em,row sep=2.1em]
&&W\arrow[dl,"a\;(\sim)"']\arrow[dr,"b"]&&\\
&U\arrow[dl,"s\;(\sim)"']\arrow[dr,"f"]&&V\arrow[dl,"t\;(\sim)"']\arrow[dr,"g"]&\\
X&&Y&&Z.
\end{tikzcd}
\]
由 \(fa=tb\)，在局部化中有
\[
 Q(t)^{-1}Q(f)=Q(b)Q(a)^{-1}.
\]
因此两个屋顶的复合为
\[
 Q(g)Q(t)^{-1}Q(f)Q(s)^{-1}
 =Q(gb)Q(sa)^{-1},
\]
正是沿图的外侧两条路径得到的屋顶 \((sa,gb)\)。图中标了 \(\sim\) 的 \(s,t,a\) 必须为拟同构，才能使用上述逆；辅助箭头 \(b\) 无须为拟同构。若另一选择为 \((a',b')\)，先共同细化拟同构 \(a,a'\)，得到 \(ac=a'c'\)。于是 \(tbc=tb'c'\)，用引理第二项再前复合一个拟同构，便使 \(bc=b'c'\) 也成立。因此两复合屋顶等价。更换原屋顶的代表时，把其共同细化与这一计算合并，同样得到等价的复合。

三重屋顶的两种括号分别使用有限个这样的交换方块。先将它们涉及的分母逐一共同细化，再用消去条件使所有指向同一对象的两条有限路径相等，就得到一个同时计算两种括号的屋顶；其左边是各分母的复合，右边是各分子的复合。因此结合律归结为 \(K(\mathcal A)\) 的结合律。这里每次只处理三重复合中的有限个方块，不涉及无限图式的同时提升。恒等屋顶为 \((1,1)\)，单位律由取恒等交换方块得到。

函子 \(Q\) 把 \(f\) 送到 \((1,f)\)。若 \(s\) 拟同构，\((s,1)\) 是 \(Q(s)\) 的逆。任意倒转拟同构的函子 \(F\) 必须把屋顶送到 \(F(f)F(s)^{-1}\)；共同细化保证此式与代表无关，复合公式保证它保持复合。这就给出唯一的分解函子 \(\overline F\)。

还需核对自然变换的分解。设 \(F,G:K(\mathcal A)\to\mathcal E\) 都把拟同构送为同构，\(\eta:F\Rightarrow G\) 为自然变换。屋顶范畴的对象不变，所以只能取 \(\overline\eta_X=\eta_X\)。对屋顶 \(X\xleftarrow{s}U\xrightarrow{f}Y\)，由 \(\eta\) 的自然性有
\[
 \eta_YF(f)=G(f)\eta_U,\qquad
 \eta_XF(s)=G(s)\eta_U.
\]
因 \(F(s),G(s)\) 可逆，得到
\[
 \eta_Y\overline F(s,f)
 =G(f)\eta_UF(s)^{-1}
 =G(f)G(s)^{-1}\eta_X
 =\overline G(s,f)\eta_X.
\]
因此这组分量确实给出唯一的自然变换 \(\overline\eta:\overline F\Rightarrow\overline G\)。屋顶范畴满足局部化的全部泛性质，正是 \(D(\mathcal A)\)。最后把两个普通态射看作分母恒等的屋顶，共同细化条件便是所写的相等判据；后复合的判据由引理第三、第四项得到。
\end{proof}

对两个屋顶 \(X\xleftarrow{s}U\xrightarrow{f}Y\)、\(X\xleftarrow{t}V\xrightarrow{g}Y\)，用通分条件取 \(a:W\to U\)、\(b:W\to V\)，使 \(sa=tb\) 为拟同构。二取三性质保证 \(a,b\) 也都是拟同构。两张右箭头此时有同一源对象，才可相加，所得屋顶为
\[
 X\xleftarrow{sa}W\xrightarrow{fa+gb}Y.
\]
同分母时，这就是 \((s,f)+(s,g)=(s,f+g)\)。不同的通分选择可进一步共同细化，并用消去条件使相应右箭头相等；预复合又保持加法，故结果与选择无关。复合公式及原来的双线性给出新复合的双线性。零对象和有限双积由 \(K(\mathcal A)\) 保留下来。因此 \(D(\mathcal A)\) 是加性范畴，\(Q\) 为加性函子。屋顶演算的一般形式可参阅 Weibel\cite[Theorem 10.3.7; Corollaries 10.3.8--10.3.11]{Weibel1994HomologicalAlgebra}。

\subsection{从同伦范畴到导出范畴}
\label{b4:subsec:1702:03}

\begin{theorem}\label{b4:thm:derived-triangulated}
设 \(\mathcal A\) 为局部小交换范畴，并采用\subsecref{b4:subsec:1701:03}的集合大小约定，记局部化函子为 \(Q:K(\mathcal A)\to D(\mathcal A)\)。在 \(D(\mathcal A)\) 中，以平移诱导的自等价为平移，以同构于 \(K(\mathcal A)\) 中好三角之像的三角为好三角，则 \(D(\mathcal A)\) 是三角范畴，\(Q\) 保持平移和好三角。
\end{theorem}
\begin{proof}
平移保持拟同构，所以经局部化诱导自等价。(TR1) 中的恒等三角和对同构封闭性已成立。任意屋顶 \(X\xleftarrow{s}U\xrightarrow{f}Y\) 的左箭头在 \(D\) 中可逆，因此 \(f\) 的锥三角运送后给出该屋顶的好三角。(TR2) 由 \(K\) 中的旋转直接得到。

核对 (TR3) 时，先把两个三角换成 \(K\) 中标准三角的像，首箭头记为 \(f:X\to Y\)、\(f':X'\to Y'\)。给定导出范畴中的交换方块 \((\alpha,\beta)\)，分别取屋顶
\(X\xleftarrow{s}S\xrightarrow{a}X'\)、\(Y\xleftarrow{t}T\xrightarrow{b}Y'\)。沿拟同构 \(t\) 通分 \(fs\)，得到拟同构 \(u:U\to S\) 和 \(k:U\to T\)，满足 \(t k=fsu\)。方块在 \(D\) 中交换，故 \(Q(f'au)=Q(bk)\)。由屋顶相等判据，再取拟同构 \(v:V\to U\)，使 \(f'auv=bkv\) 在 \(K\) 中成立。

上述构造给出 \(K\) 中的两个交换方块，它们共享上方箭头 \(kv:V\to T\)，其交换等式为
\[
 t(kv)=f(suv),\qquad b(kv)=f'(auv).
\]
把 \(kv\) 补成一条标准锥三角，再分别用 \(K\) 中的 (TR3) 映向 \(f\)、\(f'\) 的三角：前一对分量为 \((suv,t)\)，后一对为 \((auv,b)\)。这些方块的等式在 \(K\) 中成立，所选的复形映射代表可以只同伦交换；前章的锥映射公式正是用同伦修正它们。前一补全的第三个分量是拟同构，因为前两个分量是拟同构，对各次上同调长正合列使用\cref{b4:thm:five-lemma}即可。因此这一三角态射在 \(D\) 中成为同构；先取其逆，再复合后一三角态射，就得到原方块的补全。这证明 (TR3)，并使\cref{b4:lem:pretriangle-hom}现在也可在 \(D\) 中使用。

对 (TR4)，设 \(\alpha:X\to Y\)、\(\beta:Y\to Z\) 为两个导出态射。要对它们实际构造锥，需先用屋顶通分，将这一对导出态射换成 \(K\) 中可复合的复形映射。先取 \(\beta\) 的屋顶 \(Y\xleftarrow{t}T\to Z\)，再取 \(\alpha\) 的屋顶 \(X\xleftarrow{s}S\to Y\)，并沿 \(t\) 通分右边。得到 \(K\) 中一串 \(U\to T\to Z\)，以及拟同构 \(U\to X\)、\(T\to Y\)，其局部化像正好代表原来的可复合对。在 \(K\) 中对此串应用\cref{b4:thm:homotopy-triangulated}证明中的显式八面体，再施加 \(Q\) 并沿上述两个同构运送，便得到补全 \(\alpha,\beta,\beta\alpha\) 的三个好三角及它们的八面体。

这样得到的是由所选复形映射构造的三角，而 (TR4) 还要求八面体适配任意事先指定的三个好三角。对每个首箭头 \(f\in\{\alpha,\beta,\beta\alpha\}\)，在所得三角与指定三角的前两项取恒等态射，由已证的 (TR3) 补出第三分量 \(c_f:C_f\to C_f'\)。\cref{b4:lem:pretriangle-hom}只需前三条公理，故可在这里应用：前两个分量为同构，所以 \(c_f\) 也可逆。把所得八面体在第三对象之间的箭头记为 \(a:C_\alpha\to C_{\beta\alpha}\)、\(b:C_{\beta\alpha}\to C_\beta\)、\(\gamma:C_\beta\to C_\alpha[1]\)，定义运送后的箭头为
\[
 a'=c_{\beta\alpha}a c_\alpha^{-1},\qquad
 b'=c_\beta b c_{\beta\alpha}^{-1},\qquad
 \gamma'=c_\alpha[1]\gamma c_\beta^{-1}.
\]
各个 \(c_f\) 都与对应三角的后两条箭头相容，故八面体的四个相容等式保持成立；最后一个面经 \((c_\alpha,c_{\beta\alpha},c_\beta)\) 成为同构的好三角。这便满足针对任意指定三角的 (TR4)。
\end{proof}

\begin{proposition}\label{b4:prop:derived-short-exact-triangle}
设 \(\mathcal A\) 为局部小交换范畴，\(0\to A\xrightarrow{i}B\xrightarrow{q}C\to0\) 是 \(\mathcal A\) 中上链复形的短正合列，不要求逐次分裂。记 \(Q:K(\mathcal A)\to D(\mathcal A)\) 为局部化函子，\(p_i:\operatorname{Cone}(i)\to A[1]\) 为锥的投影，令 \(r:\operatorname{Cone}(i)\to C\) 为 \(r(b,a)=q(b)\)。则 \(r\) 是拟同构，且
\[
 A\xrightarrow{i}B\xrightarrow{q}C\xrightarrow{\delta}A[1],
 \qquad\delta=Q(p_i)Q(r)^{-1},
\]
是 \(D(\mathcal A)\) 中的好三角。
\end{proposition}
\begin{proof}
\(r\) 是逐次满的复形映射，因为 \(q\) 逐次满且 \(qi=0\)。利用 \(i^n\) 识别 \(\ker q^n\) 与 \(A^n\)，有
\[
 (\ker r)^n=\ker q^n\oplus A^{n+1}
 \cong A^n\oplus A^{n+1}.
\]
在此识别下，锥的微分为
\[
 \mathrm{d}_{\ker r}^n(a,a')
 =(\mathrm{d}_A^na+a',-\mathrm{d}_A^{n+1}a'),
\]
恰好是 \(\operatorname{Cone}(1_A)\) 的微分。该锥由\cref{b4:prop:cone-sign-compatibility}可缩，因而零调。对
\(0\to\ker r\to\operatorname{Cone}(i)\xrightarrow{r}C\to0\)
应用\cref{b4:thm:cohomology-long-exact}，得到每个 \(H^n(r)\) 为同构，故 \(r\) 拟同构。又因 \(r\) 与锥的包含复合等于 \(q\)，把 \(i\) 的标准锥三角经 \(Q(r)\) 运送，即得到所述好三角及最后箭头 \(Q(p_i)Q(r)^{-1}\)。这只使用 \(r\) 的拟同构性，无须原短正合列具有复形截面。
\end{proof}

若原短正合列逐次分裂，选 \(B^n=A^n\oplus C^n\)，并写
\[
 \mathrm{d}_B=\begin{pmatrix}\mathrm{d}_A&t\\0&\mathrm{d}_C\end{pmatrix}.
\]
由 \(\mathrm{d}_B^2=0\) 得 \(\mathrm{d}_At+t\mathrm{d}_C=0\)。前章\cref{b4:prop:degreewise-split-triangle}中的复形映射截面，在这里记为
\[
 \sigma^n:C^n\longrightarrow\operatorname{Cone}(i)^n,
 \qquad \sigma^n(c)=((0,c),-t^n(c)).
\]
它的微分满足
\[
 \mathrm{d}_{\operatorname{Cone}(i)}\sigma(c)
 =((0,\mathrm{d}_Cc),\mathrm{d}_At(c))
 =\sigma(\mathrm{d}_Cc),
\]
且 \(r\sigma=1_C\)。局部化后 \(Q(r)\) 可逆，因而 \(Q(\sigma)=Q(r)^{-1}\)，于是
\[
 \delta=Q(p_i)Q(r)^{-1}=Q(p_i\sigma)=Q(-t).
\]
负号来自 \(\sigma\) 的第二分量，与前章的正投影约定一致。在模的情形，普通上同调连接同态由余圈 \(c\) 的提升 \((0,c)\) 及其微分 \((t(c),0)\) 计算，给出 \([c]\mapsto[t(c)]\)；这里三角连接箭头在上同调上给出 \([c]\mapsto-[t(c)]\)。一般交换范畴中的比较已由\cref{b4:prop:degreewise-split-triangle}之后的余圈子对象计算给出，同样相差所选锥约定的负号。

\ProseParagraphBreak
此外，\(D\) 中态射为同构，当且仅当其锥为零：用三角的 Hom 正合列，锥为零使该态射对每个 Hom 群诱导同构，从而可逆；反向由三角上同调或同样的 Hom 列得到。对原来的复形映射，这恰好恢复拟同构判据。
